\ifspace
\paragraph{The Trade-off Between Security and Scalability.}
\else
\section{The Trade-off Between Security and Scalability.}
\fi
As illustrated in \autoref{tab:compare}, all prior systems fail to simultaneously achieve safety, liveness, and storage scalability.
Conventional BFT protocols and certain blockchain databases~\cite{peng2020falcondb} fully replicate their states across all nodes, ensuring safety and liveness but lacking storage scalability.
\textit{BlockchainDB}, SharPer, and \textsc{GeoBFT} necessitate pre-defined fault-tolerant clusters.
They can only tolerate a fraction of faulty nodes within each cluster and fail when the total number of faulty nodes overwhelms any single cluster.
\textsc{Elastico}, OmniLedger, and GearBox's fixed-size, randomly-sampled committees are designed to handle low global fault rates (\eg, $\frac{1}{4}$).
Increasing the fault fraction without adjusting the committee size leads to immediate safety and liveness failures.
To maintain the same security level \lijl{what properties?} \gd{Done.}, the committee size must be adaptively increased based on the actual fault fraction to ensure that committees contain no more than $\frac{1}{3}$ faulty members with high probability.
For instance, with a $\frac{1}{4}$ fault probability, sampling 378 nodes from 1000 is required to achieve a $10^{-6}$ faulty committee probability, whereas a $\frac{3}{10}$ fault probability necessitates sampling at least 798 nodes for the same probabilistic safety and liveness guarantees\footnote{These numbers are derived with hypergeometric distribution~\cite{hypergeometric-distr} as formulated in the original works.}.
As the fault probability approaches the desired $\frac{1}{3}$, the required number of sampled nodes nears the total number of nodes, and the system reverts to a fully replicated architecture, losing storage scalability.
The other systems employing randomized committees rely on additional assumptions for their safety guarantees, such as trusted hardware or synchronous network conditions.
Monoxide appears to secure all three properties concurrently.
However, if we adhere to the work's assumption that professional mining facilities mining across all zones dominate the mining power, storage scalability is lost because nodes must store states from every zone.
Conversely, if nodes engage in solo mining~\cite{bitcoin_solo_mining} and operate on single zones, individual zones become susceptible to subversion by faulty mining power, resulting in the loss of safety and liveness.
Ultimately, no prior system can simultaneously achieve all three properties.